\section{A coherent tetrahedral-boundary model}\label{sec:tetra}

Let $p_1,p_2,p_3,p_4$ be affinely independent protected source witnesses in
an open convex ambient set $X$, and put
\[
 T=\conv\{p_1,p_2,p_3,p_4\}.
\]
After a small homothety about the barycenter if necessary, the four witnesses
lie in a bounded relative carrier compactly contained in $X$.

For each neuron and each of the six edges $e_{rs}=[p_r,p_s]$, convexity gives
a relatively open interval trace, possibly empty or the whole edge.  Record
all six traces in one common affine parametrization before treating the
faces.

\begin{theorem}[Coherent tetrahedral-boundary model]\label{thm:tetra}
The four triangular faces of $T$ admit polygonal neural-code models such
that:
\begin{enumerate}
\item the code on each face equals the source code restricted to the relative
      interior of that face and is therefore a subcode of the global source
      code;
\item every source word occurring on a face survives;
\item every strict trace on each of the six edges is unchanged;
\item the two incident face models agree pointwise on every shared edge,
      including endpoint positions and simultaneous events.
\end{enumerate}
Thus $\partial T$ carries a coherent finite facewise polygonal output model.
\end{theorem}

\begin{proof}
Apply \cref{thm:main} to each triangular face in its relative affine plane,
prescribing its three edge-supporting lines and protecting the desired face
witnesses.  A shared edge carries one source subset for every neuron, and
both incident applications preserve it exactly.  The statements follow.
\end{proof}

\begin{figure}[t]
\centering
\begin{tikzpicture}[scale=.95]
  \coordinate (p1) at (0,0);
  \coordinate (p2) at (4.4,0);
  \coordinate (p3) at (1.3,2.7);
  \coordinate (p4) at (2.15,1.15);
  \fill[gray!8] (p1)--(p2)--(p3)--cycle;
  \draw[thick] (p1)--(p2)--(p3)--cycle;
  \draw[thick] (p1)--(p4)--(p2);
  \draw[thick] (p3)--(p4);
  \draw[dashed] (p4)--(p1);
  \foreach \x/\y in {p1/1,p2/2,p3/3,p4/4}{\fill (\x) circle (2pt) node[above right] {$p_\y$};}
  \draw[blue!70!black,line width=2.7pt] ($(p1)!.18!(p2)$)--($(p1)!.65!(p2)$);
  \draw[orange!85!black,line width=2.7pt] ($(p1)!.35!(p3)$)--($(p1)!.82!(p3)$);
  \draw[green!50!black,line width=2.7pt] ($(p2)!.2!(p3)$)--($(p2)!.72!(p3)$);
  \draw[purple!70!black,line width=2.7pt] ($(p1)!.25!(p4)$)--($(p1)!.78!(p4)$);
  \draw[red!70!black,line width=2.7pt] ($(p2)!.32!(p4)$)--($(p2)!.88!(p4)$);
  \draw[cyan!60!black,line width=2.7pt] ($(p3)!.18!(p4)$)--($(p3)!.7!(p4)$);
  \node at (5.05,1.75) {six prescribed edge traces};
\end{tikzpicture}

\caption{All six source edge traces are frozen first; the four faces are then
polygonalized independently but coherently along the shared edges.}
\label{fig:tetrahedron}
\end{figure}

\begin{remark}[What is not supplied]
The theorem does not construct convex three-dimensional polytopes whose
boundary traces are the four face models.  In particular, it does not match
numerical gauge values, supporting facets, or inward normal derivatives
across an edge.  Those are genuinely additional three-dimensional
compatibility conditions.
\end{remark}
